\section{What happened during the run}
\label{sec:trajectory}

The chronology has two complementary sources. Commits give dated,
recoverable states of the repository. The exported messages give a
selective account of what the agent believed it had established at each
stage. Neither records the instant at which an idea first occurred.
We therefore describe observable transitions---a proof being saved, a
certificate being completed, a claim being withdrawn---rather than
assigning precise discovery times to unobserved reasoning.

\subsection{Selection, short arguments, and early corrections}

The initial plan favored finite counterexamples, explicit constructions,
linear algebra, and questions that could admit short certificates.
This was an adaptive search policy over the whole Notebook. It helps
explain the eventual collection: the agent repeatedly moved to another
entry when a search became expensive or the required generalization
remained unclear. A problem's age was not a measure of the time spent on
it in this experiment, and the initial index was not a promise to attempt
every entry.

The first retained mathematical lead was \kn{21.106}. Message~95
announced a possible formula selecting generators of the infinite cyclic
center of the integral Heisenberg group; message~152 reported a written
candidate and local commit. Section~\ref{subsec:concise} gives the
argument. At the same time the agent launched finite screens for
other questions, including transitive groups, character tables, and
positive binary words. This pattern continued: short symbolic tasks
occupied the foreground while finite searches ran in separate processes.

The early graph-symmetry search exposed a substantive validation error.
For \kn{21.52}, normalizing a permutation representation of inner
automorphisms did not by itself certify that a permutation of involutions
came from an automorphism of the original group. Message~653 records
the replacement of that condition by construction of an actual group
automorphism and verification of its action on every involution.
The lesson concerns the mathematical predicate being checked: a correct
program for a weaker predicate can still give a misleading answer to the
research question.

Literature searches also changed the accounting. A graded nil-algebra
construction for \kn{21.132} initially appeared on the candidate list,
but was later excluded as prior work. For \kn{13.19}, a small explicit
quotient construction passed exact checks before a 2018 result was
confirmed to imply the negative answer. Message~1155 explicitly
removed it from the new-solution count while retaining the construction.
The longer episode around \kn{17.34} is particularly informative.
The agent wrote a PBW-based argument and ran thirteen exact controls;
it first found prior coverage in nilpotency class two, then found a
strong-amalgamation theorem with the full required scope. Messages~2729
and~2859 document the intervening increase and subsequent decrease in
the working candidate count. These are rediscoveries or alternative
deductions with potential expository value, but they cannot support a
claim of new priority.

\subsection{Constructions guided by the literature}

Two examples in Section~\ref{sec:examples} illustrate how source reading
and explicit calculation interacted. For \kn{21.68}, the order-96
configuration in Kida's work suggested an extension by an abelian
group that could preserve a nonmonomial character. The resulting
group has order \(2592\). Its membership in the class of semi-abelian
groups and its failure to be monomial require different arguments;
checking only one would not settle the conjecture.

For \kn{20.108}, exceptions in a multiple-holomorph theorem suggested
a group of order \(605\). The retained proof describes its entire
automorphism group, its holomorph as explicit permutations, and a
normalizing permutation whose coset has order four. The source supplied
a productive place to look; the final certificate does not ask a reader
to trust a SmallGroups identifier or an unexplained normalizer
calculation. In both examples the finite checks are useful controls on
conventions and coordinates, while the displayed construction explains
why the claimed property holds.

Other entries depended more substantially on imported theorems.
The constructions for Problems~10.62 and~4.75 share the same geometric
existence theorem, and the carpet arguments for Problems~19.61 and~19.62
share structural machinery. The inventory retains separate Notebook
questions, but they are not independent pieces of evidence about
discovery ability. The appendices identify these dependencies.

\subsection{A long computation and the meaning of completion}

Problem~20.100 became the largest verification task. A structural
reduction transformed a statement about arbitrary groups into an exact
finite problem for each fixed value of \(n\). Small values could then be
established by certificates of integer relations, without imposing a
bound on the order of the original group. Messages~6458--6763 distinguish
the reduction, prototype searches, and completed checks for
\(n=4,5,6\).

For \(n=7\), generation produced a certificate with \(42{,}891{,}332\)
nodes. Message~10686 reports its completion and immediately notes that
verification remained outstanding. Six GAP processes subsequently
checked the separate parts. The final records account for
\(786{,}577{,}344\) collision cases and \(5{,}447{,}182{,}055\)
relation implications; the complete compressed certificate occupies
\(1{,}271{,}256{,}410\) bytes. The retained full verification took
approximately 26 hours of elapsed time. A parser restricted to the
certificate's data grammar was also used, avoiding reliance on
executing arbitrary certificate text as GAP code.

This distinction between generation and checking matters. Intermediate
messages report how many nodes had been checked, how many workers were
still running, and which shards had completed. They do not convert
partial progress into a full theorem. At the deadline the fixed cases
\(n=4,5,6,7\) were recorded as established for all eligible groups.
The general problem for arbitrary \(n\) remained open and contributed
zero entries to the complete-candidate count.

Several other large searches ended with bounded information. The
endomorphism-count investigation of \kn{19.20} checked \(91{,}774\)
nonabelian groups of order at most \(511\), found no equality example,
and found \(469\) cases contradicting a stronger inequality. The
character-table screen for \kn{21.113} involved \(9{,}850\) cases:
\(7{,}573\) had the needed modular tables and \(2{,}277\) did not.
For \kn{21.99}, all \(4{,}722\) transitive groups of degrees \(2\)
through \(20\) in the installed catalogue were screened, with a separate
direct control on \(86\) groups in degrees \(2\) through \(8\).
The word search associated with \kn{18.43} processed
\(3{,}933{,}931{,}043\) positive binary necklaces of lengths \(1\)
through \(36\) without producing an identity candidate. These are
coverage statements. They do not establish the unrestricted assertions.

\subsection{Late structural results and deadline review}

The later work included structural progress on almost Engel elements
and \(p\)-Jordan bounds. For Problems~20.89/19.83 the deadline result
covers all linear groups in characteristic zero and several further
classes; the unrestricted positive-characteristic question remains
unresolved. For \kn{21.121}(b), the archive contains bounds for
soluble finite subgroups, examples with exactly computed exponents,
and permanence results. The universal proposed bound of three remains
unproved. These partials illustrate why a single binary solved/unsolved
label would discard much of the experiment's mathematical output.

Statement checking remained necessary even when a construction was
correct. For \kn{20.122}, an early treatment of inclusion-minimal
intersections did not settle the question about intersections of
minimum order. The final construction addresses the latter condition.
For \kn{15.76}(b), the retained claim concerns the full varieties of
fixed derived length, with a material distinction from arbitrary proper
subvarieties. The fixed representation degree in \kn{10.35} and the
minimal-polynomial convention in \kn{15.65} are similarly essential
parts of the claims.

The final review separated statement coverage from deeper proof review.
All 46 candidate entries received a fresh comparison with the supplied
PDF; 22 selected candidate proofs received additional late rereading.
This is not evidence that all 46 underwent a new, exhaustive
mathematical and literature audit in the final phase. The closeout audit
also checked \(4{,}231\) file bindings, covering \(2{,}242\) distinct
files totaling \(2{,}008{,}448{,}358\) bytes, and 92 local links.
These integrity checks establish that the named artifacts existed and
matched recorded digests. They do not establish the truth or novelty of
their contents. The research deadline and the subsequent administrative
closeout are kept separate in the chronology.

